

\chapter{Methodology}
\label{ch:method}

\section{Research Question}

This project addresses one question: \textbf{Can a CNN, trained on rendered ECG images, detect cardiac abnormalities and rhythm disorders in real world images, and how can information and confidence be made reliable and understandable to users?}
It is broken into the objectives O1--O6 of Chapter~\ref{ch:intro}, and each method
below is linked to the objective it serves.

\section{Users and Requirements}

The main users are university students learning ECG interpretation, who need to
understand why a trace was flagged. A clinician view, which would show a short list of
flags with measured values and no teaching text, is a secondary target. The
requirements were drawn from the literature on how clinicians and users respond to
explained AI output \cite{hughes-noehrer_contrasting_2024,cabitza_painting_2023} and
from the gaps identified in Chapter~\ref{ch:lit}, and are prioritised with MoSCoW
analysis (Appendix~\ref{app:moscow}). The must-have requirements are:
upload of an ECG image; display of flagged conditions with a calibrated confidence for
each; a Grad-CAM heatmap over the original image; and a clear statement that the output
supports review and is not a diagnosis. [TODO: confirm and trim to your final must-have list.]






